\section{Methods}
\label{sec:methods}

The model is a discrete hazard estimator. Each request contributes one binary
observation per bin it survives, labelled by whether its reuse distance
resolves in that bin. A gradient-boosted tree ensemble \cite{friedman2001}
predicts the conditional resolution probability from request features.

\subsection{Discrete Hazard with Gradient-Boosted Trees}

For a request with reuse distance $r_i$ observed in bin $k_i$, the expanded
data set contains rows $(i, k)$ for $k = 0, \ldots, k_i$, with label
\begin{equation}\label{eq:label}
    y_{i,k} \;=\; \mathbb{1}\{r_i = k\},
\end{equation}
and row $(i,k)$ is in the risk set only if $r_i \ge k$. The model estimates
$h_k(z_i) = \Pr(y_{i,k} = 1 \mid r_i \ge k, z_i)$, and the log-likelihood of the
observed data is
\begin{equation}\label{eq:likelihood}
    \ell \;=\; \sum_{i} \Biggl[
        \sum_{k=0}^{k_i - 1} \log\bigl(1 - h_k(z_i)\bigr)
        \;+\; \delta_i \log h_{k_i}(z_i)
    \Biggr],
\end{equation}
where $\delta_i$ is the event indicator of \cref{tab:notation}. The first term
is the survival contribution of bins the request passed through, the second is
the resolution contribution, present only when the distance is observed.

\subsection{Two Parametrizations}

\textbf{R-as-feature.} A single ensemble is trained on the expanded rows with
the bin index $k$ appended to $z_i$ as a numeric feature. The model learns one
response surface over features and bin index, which shares strength across bins
and keeps the number of fitted models at one. The risk is that the ensemble must
represent a monotone structure in $k$ implicitly.

\textbf{Per-bin.} One binary ensemble is fitted per bin $k$ on that bin's risk
set, using $z_i$ without the bin index. Each model is small and the hazard in
each bin is unconstrained, at the cost of $K+1$ fits and no sharing between
adjacent bins. Both parametrizations produce a hazard sequence that is
normalized to $[0,1]$ and, if needed, projected to be non-increasing in $k$ for
$S(k)$.

\subsection{Training Labels and Censoring Correction}

Labels follow \cref{eq:label} directly for observed distances. A censored
request with censoring time $\tau_i$ contributes rows for $k < \tau_i$ with
label $0$ and no row at $\tau_i$; it is never labelled as a resolution. Two
corrections are compared for the bias that remains when censoring correlates
with features:
\begin{enumerate}[label=(\roman*),nosep]
    \item \textbf{Risk-set restriction:} drop bins whose at-risk count falls
    below a floor, and report the curve only up to the largest well-populated
    capacity.
    \item \textbf{Inverse-censoring weights:} weight each uncensored row by the
    inverse of the estimated censoring survival function at its bin, following
    the usual inverse-probability-of-censoring construction.
\end{enumerate}
The comparison is reported in \cref{sec:experiments}; neither correction is
assumed to dominate.

\subsection{Hyperparameter Search}

Hyperparameters are selected with Optuna \cite{akiba2019} using the tree
Parzen estimator over the space in \cref{tab:searchspace}. The objective is the
validation integrated Brier score defined in \cref{sec:metrics}, not the
training likelihood, because likelihood favors over-confident hazards.

\begin{table}[ht]
\centering
\caption{Optuna search space. Values are placeholders until the sweep is run.}
\label{tab:searchspace}
\begin{tabular}{@{}l l l@{}}
\toprule
Hyperparameter & Range & Scale \\
\midrule
Number of leaves & $[15, 255]$ & integer, log \\
Learning rate & $[10^{-3}, 10^{-1}]$ & log \\
Minimum child samples & $[20, 500]$ & integer, log \\
Feature fraction & $[0.5, 1.0]$ & linear \\
Bagging fraction & $[0.5, 1.0]$ & linear \\
L2 regularization & $[10^{-8}, 10^{1}]$ & log \\
Max bins & $\{63, 127, 255\}$ & categorical \\
Hazard bins $K$ & $\{32, 64, 128\}$ & categorical \\
\bottomrule
\end{tabular}
\end{table}

\subsection{Evaluation Metrics}
\label{sec:metrics}

Let $\hat{H}$ be the estimated curve and $H$ the simulated curve on a held-out
segment of the same trace, evaluated on the capacity grid $\mathcal{C}$.
\begin{enumerate}[label=(\roman*),nosep]
    \item \textbf{MRC error:} the maximum absolute deviation
    $\max_{C \in \mathcal{C}} |\hat{H}(C) - H(C)|$, and the mean absolute
    deviation over $\mathcal{C}$.
    \item \textbf{Calibration:} reliability of the hazard predictions, binned by
    predicted probability, summarized as expected calibration error.
    \item \textbf{Brier score:} mean squared error of the binary hazard
    predictions, and the integrated Brier score over bins, which aggregates the
    score across the survival path rather than at one horizon.
\end{enumerate}

\subsection{Simulation Ground Truth}

Ground truth is a full simulation of the trace on the capacity grid, using
exact stack distance where the trace fits in memory and SHARDS
\cite{waldspurger2015} with a fixed sampling rate where it does not. The
sampling rate and the resulting confidence band are recorded with each curve,
so that model error is never reported below simulation error. Ground truth is
computed once per trace and capacity, never per model configuration.
